
% \label{tab:atributos}
% \vspace{-0.3cm}
% \end{table}

%\ligia{corta as tabela e menciona no git.}

% \begin{table}[htbp]
% \centering
% \caption{Atributos Evaluated by AUTO-SEE technique} 
% \begin{tabular}{|c|c|}
% \hline
% \textbf{Atributo}  \\ \hline
% %eth - enlace
% 1-2. Total de endereços MAC de origem/destino únicos\\ \hline 
% 3. Versão mais frequente do protocolo que estará na carga útil \\ \hline 
% %ip-rede
% 4-5. Tamanho do menor/maior pacote \\ \hline 
% 6. Soma do tamanho dos pacotes \\ \hline


% 7-9. Total de pacotes do tipo ICMP/UDP/TCP\\ \hline 
% \textbf{10-11. Total de endereços IP de origem/destino únicos}\\ \hline 
% 12-16. Maior/Menor/Soma/STD/Mean do time to live dos pacotes\\ \hline 
% \textbf{17. Total de pacotes trafegados na rede} \\ \hline 

% %tcp/udp-transporte
% 18-19. Porta de origem/destino mais comuns encontrada nos pacotes\\ \hline 
% 20. Total de flags TCP únicas \\ \hline 
% 21-26. Total de flags TCP FIN/SYN/RST/PUSH/ACK/URG\\ \hline 

% 27-31. Maior/Menor/Soma/STD/Mean do valor do campo TCP window size \\ \hline 
% 32-36. Maior/Menor/Soma/STD/Mean do valor do campo TCP Sequence Number\\ \hline 

% %tempo de chegada dos pacotes
% 37-41. Maior/Menor/Soma/STD/Mean do tempo de chegada entre pacotes\\ \hline 
% 42-46. Maior/Menor/Soma/STD/Mean do tempo de chegada entre pacotes TCP \\ \hline 
% 47-51. Maior/Menor/Soma/STD/Mean do tempo de chegada desde o primeiro pacote TCP \\ \hline 


% \end{tabular}
